\section{Introduction}
\label{sec:intro}

A diffusion model generates data by reversing a noising process, and the only object it has to learn is a denoiser:
for every noise level $t$, the best guess of the clean data $a$ given a noisy version $x$
\cite{sohl2015deep,ho2020denoising,song2021scorebased}. By Tweedie's formula the score that drives the reverse
process is an affine function of the posterior mean $\E[a\mid x]$ \cite{robbins1956empirical,efron2011tweedie}, so
learning the score and learning the denoiser are the same problem. Neural networks solve it from samples, at every
$t$ at once, and what they learn about the structure of the data is an active question
\cite{kadkhodaie2024generalization,kamb2024analytic,biroli2024dynamical}.

We study this question on data whose structure is known exactly: first-order Markov sequences. A Markov sequence is
fixed by its transition kernel $K(a_{k+1}\mid a_k)$. This covers a wide class of data with a direction or a
locality, from time series of causal processes to one-dimensional physical fields with nearest-neighbour interactions,
whose stationary law is a chain. The question is what a denoiser gains from knowing that the data are Markov, and how
it should use that knowledge.

The answer has three parts, which organise the paper.
First, the structure is \emph{not} in the score. For the Gaussian chain the score is linear and its matrix is
tridiagonal at $t=0$, but diffusion fills it in: at intermediate noise every site depends on every other
(\cref{sec:gaussian}). A denoiser constrained to be local, by a tridiagonal fit or a small receptive field, has an
error floor that no amount of data removes.
Second, the structure \emph{is} in the posterior. A Markov chain observed through independent noise is a tree, so
belief propagation (BP) computes $\E[a\mid x]$ exactly, for any kernel and at every $t$ (\cref{sec:tree}).
Third, learning the denoiser therefore reduces to learning the kernel. We do this by expectation--maximisation (EM),
with BP as the E-step, from noisy sequences only (\cref{sec:laplace}). BP enters the learning only through expected
transition counts, so the kernel can be any differentiable family, a small neural network included, without
differentiating through BP.

We test this on a non-Gaussian Markov chain with Laplace innovations, the simplest case where the exact score is not
linear, against six neural denoisers trained on the same data, and we measure how the comparison changes as the data
depart from the Markov property (\cref{sec:laplace}). We then apply it to the stochastic Allen--Cahn field, where the
choice of kernel family matters (\cref{sec:allen-cahn}).

\note{Related work in one paragraph here or in Sec.~7: HMM and Kalman EM \cite{baum1970,shumway1982em}; diffusion
from corrupted data with EM \cite{daras2023ambient,rozet2024learning,bai2024em}; BP and transformers on structured
data \cite{garnierbrun2024transformers}; analytic denoisers \cite{kamb2024analytic}. J\'er\^ome will help with the
literature.}
